\documentclass[11pt]{article}
\usepackage{mathblatt}
\newcommand{\bereich}[1]{\par\addvspace{14pt}\noindent{\large\bfseries #1}\par\nobreak\vspace{4pt}}
\newcommand{\bl}[1]{\par\noindent #1\par\vspace{3pt}}
\newcommand{\blrand}[1]{\par\noindent{\color{mbgrau}#1}\par\vspace{3pt}}
\newcommand{\blhier}[1]{\par\noindent\llap{$\blacktriangleright$\hspace{0.5em}}\textbf{#1}\par\vspace{3pt}}
\newcommand{\stufe}[1]{\par\noindent\hspace*{1.5em}{\small #1}\par\vspace{2pt}}
\newcommand{\pkt}{\textperiodcentered{}}
\begin{document}

\pfheftstil
\pflaufkopf[P10 \pkt{} Prüfung]{Quadratische}{PA40-4}\pfseitevon
\pfheftkopf{Quadratische}{}
\bl{Parabel}
\blhier{Quadratische Gleichung}
\stufe{Nullstellen gesucht}
\stufe{x-Wert gesucht}
\stufe{Schnittpunkte gesucht}
\stufe{Gerade und Parabel}

\par\addvspace{14pt}\noindent Vorher: Lineare\par
\par\noindent Weiter: Gleichungssysteme\par
\end{document}